\chapter{Our setup}
\label{ch:experimental-setup}

What did the FLFM cameras record, and which of those measurements can be used
to assess the reconstruction in Chapter~\ref{ch:image-reconstruction}? Answering that requires the optical
layout, the acquisition record, and the processing run that produced each
result. This chapter describes the instrument and uses the June~3, 2026
eyedrop acquisition as a case study for the data and comparison protocol.

\section{Optical configuration}
\label{sec:optical-configuration}

The instrument has one illumination path and three camera paths. The main
Fourier light-field microscopy (FLFM) camera supplies the data for reconstruction. ThorCam supplies a
conventional image for comparison, while the opposite-side Validation Camera
provides another view of the sample. A Coherent
SAPPHIRE 488-200 CDRH laser supplies $488\,\mathrm{nm}$ light. An expanding
lens and a $40\times$ objective focus it onto the sample. That objective also
collects light for the main camera or, when the flip mirror is engaged, the
ThorCam. A second $40\times$ objective faces the sample from the other side
and leads to the Validation Camera. The component inventory gives
$20\,\mathrm{mW}$ as a commonly used laser setting. It does not give the
power at the sample or the setting for individual acquisitions.

In the main collection path, a two-lens relay images the objective's Fourier
(pupil) plane onto the microlens array. The installed Advanced Microoptics
Systems GmbH APO-Q-P3000-R23.5 has $10\times10$ fused-silica microlenses on
a square grid with $3\,\mathrm{mm}$ pitch. Its working focal length is
$55\,\mathrm{mm}$; the flat side faces away from the incoming collimated
light. Another two-lens relay takes the array output to the Teledyne-FLIR
Blackfly S BFS-PGE-244S8M main camera. In the ray-optics design, the lenses
before the array have focal lengths of $250\,\mathrm{mm}$ and
$750\,\mathrm{mm}$, and those after it have focal lengths of
$400\,\mathrm{mm}$ and $150\,\mathrm{mm}$.
Individual lens part numbers and measured separations are still missing.

Before the microlens array, a flip-mounted mirror can send the collected light
to the ThorCam branch. This branch avoids the array and focuses through a
regular lens onto a Thorlabs CS2100M-USB camera. The ray-optics script models
that lens as $300\,\mathrm{mm}$. The objective and first relay are shared
with the FLFM path; everything after the mirror belongs to one branch or the
other. Switching the mirror means that the two cameras capture at different
times. Across the sample, the validation path passes through its own
$40\times$ objective and a lens modelled as $100\,\mathrm{mm}$ before
reaching a Teledyne-FLIR Blackfly S BFS-U3-122S6M camera.

The objective used for both the illumination/collection path and the validation camera path is ThorLabs N40X-PF $40\times$ objective: focal length
$5\,\mathrm{mm}$, numerical aperture 0.75, working distance
$0.66\,\mathrm{mm}$, and entry pupil $7.5\,\mathrm{mm}$. The
available records also leave open how excitation enters the shared objective,
which filters are installed, and what contrast the Validation Camera records.
The ray-optics scripts establish the path order, but they cannot settle these
hardware details or the spacing of the assembled lenses.

\section{Samples, acquisitions, and view extraction}
\label{sec:sample-acquisition}

The June~3, 2026 eyedrop specimen used Thermo Fisher fluorescent polystyrene
microspheres of nominal $2\,\mathrm{\mu m}$ diameter. The bead suspension was
diluted at a reported $1{:}10$ volume ratio with double-distilled water.
An aliquot of $30\,\mathrm{\mu L}$ was applied to a thin microscope slide and
air-dried. One drop of Systane gel eyedrops was then placed directly over the
dried beads on the same slide and remained there during imaging. The dilution
convention (stock-to-water or stock-to-final-volume), bead concentration,
drop volume and thickness, slide thickness, and any additional coverglass are
not recorded. The position of this sample relative to each objective also
needs to be established for an optical schematic.

Acquisitions are grouped by date, with the archive organization summarized in
Table~\ref{tab:acquisition-run-archive}. Folder proximity alone does not
establish that two files depict the same capture; scene identity and
acquisition date are essential for matching the sequential FLFM and ThorCam
images. The June~3 case has an FLFM
capture labelled \texttt{beads\_2um\_dense4\_eyedrop\_4kus}, with ThorCam and
Validation Camera files bearing the same scene label. Their capture timing,
exposure and gain settings, and image correspondence remain unconfirmed in
the available acquisition records.

Before reconstruction, the FLFM sensor image is cut into a square grid of
elemental views. Write a view as $V_{ij}(y,x)$, where
$i,j\in\{0,\ldots,N-1\}$ index the grid and $(y,x)$ are pixel coordinates
within a view. The resulting array has shape $(N,N,H,W)$; a selected exposure
stack adds an exposure axis. The physical $10\times10$ array alone does not
fix the extracted grid for any particular run. Reproducing view extraction
requires the grid center, pitch or patch geometry, orientation, and retained-view
set, together with the preprocessing state of the supplied image. An unmodified
sensor
image, a dark-subtracted image, and a display-normalized image carry different
values. The archive contains dark-frame and gain-calibration files, but their
presence does not establish their use for a dataset. The same applies to
exposure fusion:
it belongs to a run only when its saved configuration enables it.

\section{Reconstruction and selection records}
\label{sec:reconstruction-records}

Chapter~\ref{ch:image-reconstruction} describes one reconstruction sequence: global translation
alignment, then local pixelwise refinement and averaging. Tracing its
intermediate global-stage average and final rectified image to an acquisition
requires the source capture, extracted view grid, processing configuration,
and saved numerical outputs.

The configuration gives requested settings, while effective values also depend
on the implementation and input geometry. The workspace and numerical view
stack show which data were averaged. Preview images help with visual
inspection, but may have a different intensity scale. Traceability also
requires the code revision, effective parameters, preprocessing, and both
intermediate and final outputs. Some older products are missing a source
revision or a complete history of figure preparation.
A place on a sweep leaderboard cannot fill those gaps.

Parameter sweeps and image-quality scores narrow the set of images worth
inspecting. A ranking applies only to the candidates that completed for that
dataset. Sharpness or contrast alone cannot tell us whether the image
reproduces the bead scene faithfully. Interpreting a selected example requires
its candidate set, ranking rule, failed or missing runs, and any manual choice.
If an external-camera image or bead annotation helped
select the candidate, comparing it with that same reference is part of the
selection process, not an independent test.

Table~\ref{tab:acquisition-run-archive} summarizes the archive organization;
the listed records are not necessarily all present for the June~3 candidate.

\begin{table}[htbp]
    \centering
    \caption{Acquisition and processing records in the archive. File and folder
    names identify record types, not confirmed correspondence between captures.}
    \label{tab:acquisition-run-archive}
    \begin{tabular}{@{}>{\raggedright\arraybackslash}p{0.29\textwidth}>{\raggedright\arraybackslash}p{0.65\textwidth}@{}}
        \hline
        Record type & Location or content \\
        \hline
        Camera captures & Date folders can contain \texttt{main} (FLFM),
        \texttt{thorcam}, and \texttt{validation} or \texttt{val}. \\
        Extracted views and run products & \texttt{algo\_run} contains extracted
        views, rectification runs, and analysis products. \\
        Sweep and manifest identifiers & Sweep specifications, experiment
        manifests, and generation and candidate identifiers describe the
        candidate search and run identities. \\
        Configuration and numerical outputs & \texttt{config.json} records
        requested settings; \texttt{workspace.npz} and numerical view stacks
        preserve processing data and outputs. \\
        Metric records and previews & \texttt{metric\_record.json} stores scores;
        preview images support visual inspection. \\
        \hline
    \end{tabular}
\end{table}

\section{Evaluation protocol}
\label{sec:evaluation-protocol}

The most direct comparison stays within one FLFM run: the average after global
alignment against the final rectified image. They share the acquisition and
parameter record, so the comparison shows the contribution of local
refinement. Visual interpretation depends on the crop and display scale.
Numerical comparison requires the saved images, a defined metric and direction,
and the same stated region in both images. Local remapping can leave zero-filled
boundaries; including different amounts of those boundaries would change the
score for reasons unrelated to the refinement itself.

External-camera comparisons depend on correspondence between acquisitions
and registration to a common image region. ThorCam and Validation Camera
images may differ from the FLFM reconstruction in pixel sampling,
orientation, field of view, and intensity scale. Resampling, geometric
registration, intensity handling, and the common crop or validity mask
therefore form part of the comparison record. Evaluation registration may use
scale and
rotation; the global reconstruction stage itself estimates translations only.
Saved validation workspaces may contain transformations and overlays, but the
registration and metric must be read from the selected comparison rather than
inferred from current software defaults. Sample change between the sequential
FLFM and ThorCam captures is another limit on interpretation.

The optical paths and archive structure are known well enough to locate the
June~3 candidate comparison. Its sample preparation is documented above;
exposure and power settings, extraction geometry, camera registration, and
evaluation metrics remain to be established. Any comparison based on this
single acquisition is a case study, not an estimate of performance across
independent samples.
